\documentclass{report}
\usepackage{amsfonts, amsmath}
\newcommand\R{{\mathbb R}}
\newcommand\Q{{\mathbb Q}}
\newcommand\N{{\mathbb N}}
\pagestyle{empty}
\begin{document}
\large
\centerline{\Huge Analisi Matematica II}
\centerline{\large Corso di Laurea di Informatica}
\renewcommand{\thefootnote}{ } 
\footnote{\sl Avete 2.30 ore di tempo. Ogni esercizio
vale nove punti. La sufficienza si ottiene con un punteggio $\geq$
18. Solo le risposte chiaramente giustificate saranno prese in
considerazione.}
\renewcommand{\thefootnote}{\arabic{footnote}} 
\vskip.1cm
\centerline{\Large Appello 7-07-2003}
\vskip1cm
\begin{enumerate}
\item Si determini il raggio di convergenza della serie di potenze
\[
\sum_{n=0}^\infty n^{\ln n}x^n.
\]
\item Si risolva il problema di Cauchy
\[
y''=-2y'
\]
\[
y(0)=0\;;\;\;y'(0)=1.
\]
\item Quale \`e il numero maggiore tra $(1000!)^{\frac 1{1000}}$ \ e
\ $500$\  ?
\item Calcolare la serie di Fourier della funzione $e^{-|x|}$
sull'intervallo $[-\pi,\pi]$. Sia $g:\R\to\R$ la funzione definita
dalla serie di Fourier, si dica se $g$ \`e ovunque differenziabile.
\end{enumerate}
\end{document}
